%% Согласно ГОСТ Р 7.0.11-2011:
%% 5.3.3 В заключении диссертации излагают итоги выполненного исследования, рекомендации, перспективы дальнейшей разработки темы.
%% 9.2.3 В заключении автореферата диссертации излагают итоги данного исследования, рекомендации и перспективы дальнейшей разработки темы.
\begin{enumerate}[beginpenalty=10000, midpenalty=10000, wide=\parindent, leftmargin=0pt]
    \item Проведен системный анализ методов автоматизации рабочих процессов наземных транспортно-технологических машин, на основе которого сформулирована формальная постановка задачи многозадачной координации комплекса роботизированных НТТМ. Обоснована целесообразность применения децентрализованного подхода на основе теории игр, позволяющего каждой машине автономно принимать решения и минимизировать объем обмена служебной информацией.

    \item Разработана оригинальная имитационная модель многозадачной координации, отличающаяся от известных учетом динамики изменения энергетического ресурса каждой машины $ E_i $, сложности задач $ \Theta_{T_j} $ и их временных приоритетов. Модель формализована как байесовская координационная игра на связном графе взаимодействий $ G = ([N],\mathcal{E}) $, для которой доказано существование равновесий Нэша, зависящих от структуры взаимодействия и соотношения сложностей задач.

    \item Впервые предложены две новые стратегии децентрализованного выбора задач (\textit{энергосберегающая стратегия} $ \Gamma_{\min\text{Э}} $ и \textit{стратегия минимизации временных задержек} $ \Gamma_{\min\text{ВЗ}} $) и по ним разработаны соответствующие алгоритмы.

    \item Создан программный комплекс на языке Go, реализующий предложенные стратегии, позволяющий проводить многовариантное моделирование работы комплекса РНТТМ в контролируемых условиях. В ходе вычислительных экспериментов на двух конфигурациях в трех сценариях ресурсного обеспечения установлено:
    \begin{itemize}[beginpenalty=10000, midpenalty=10000, wide=\parindent, leftmargin=0pt]
        \item энергосберегающая стратегия $ \Gamma_{\min\text{Э}} $ обеспечивает наилучшее соотношение <<выполненные задачи / потраченная энергия>> во всех сценариях;
        \item в условиях жесткого дефицита энергии $ \Gamma_{\min\text{Э}} $ выполняет больше задач для разнородных комплексов, чем лучший из аналогов;
        \item анализ коэффициента Джини подтвердил, что стратегия $ \Gamma_{\min\text{Э}} $ обеспечивает не только высокую производительность, но и статистически справедливое распределение нагрузки между машинами.
    \end{itemize}

    Согласно актам внедрения, применение разработанных алгоритмов обеспечило:
    \begin{itemize}[beginpenalty=10000, midpenalty=10000, wide=\parindent, leftmargin=0pt]
        \item рост производительности комплекса РНТТМ на $ 11{,}2 $-$ 13{,}1\,\% $;
        \item сокращение общего времени выполнения операций на $ 8{,}7 $-$ 9{,}1\,\% $;
        \item повышение энергоэффективности (снижение удельных энергозатрат) на $ 11{,}8 $-$ 12{,}3\,\% $;
        \item уменьшение времени простоев оборудования на $ 8{,}2 $-$ 9{,}1\,\% $.
    \end{itemize}

    \item Выполнена комплексная оценка технико-экономической эффективности внедрения разработанных алгоритмов. Для типового комплекса из 10 машин годовой экономический эффект за счет сокращения времени выполнения операций составляет $ 2{,}67 $ млн руб. при сроке окупаемости капитальных затрат $ 1{,}7 $ года. Результаты работы апробированы в АО <<Петербург-Дорсервис>> и парках машин Министерства обороны Российской Федерации, что подтверждено соответствующими актами.
\end{enumerate}

Таким образом, решена актуальная научно-техническая задача повышения эффективности комплекса роботизированных наземных транспортно-технологических машин за счет разработки имитационной модели и алгоритмов децентрализованной многозадачной координации, что полностью соответствует поставленной цели и подтверждает выдвинутую гипотезу исследования.
